\scanpage{2}
\begin{center}
INSTITUT\\
DES HAUTES ÉTUDES SCIENTIFIQUES\\
\rule{1.2cm}{.4pt}
\end{center}
8.4.4.

\uline{NB}

 a) $F$ opérant sur $H^i(X)\otimes_W k$ pour $i>n$ est \uline{nilpotent}, donc aussi opérant sur $\operatorname{Tors}^i(X)\otimes_W k$ ou sur $\operatorname{Tor}_{\unsure{1}}^W(H^i(X),k)\simeq{}_p\operatorname{Tors}^i(X)$.

On a $H^n(X,\underline{O}_X)_{\mathrm{ss}}\unsure{\xleftarrow{\sim}|\simeq}(H^n(X)\otimes_W k)_{\mathrm{ss}}$ d'où
\diagram{\begin{tikzcd}[column sep=small]
0\arrow[r]& (\unsure{\operatorname{Tors}|H}^{n}_{\unsure{\mathrm{ss}}}(X)\otimes k)_{\mathrm{ss}}\arrow[r]&H^n(X,\underline{O}_X)_{\mathrm{ss}}\arrow[r]& (H^n_{\mathrm{lib}}(X)\otimes_{\unsure{\mathrm{ss}}}k)_{\mathrm{ss}}\arrow[r]&0
\end{tikzcd}}
\margin{$(*)$ \unsure{$T_{\mathrm{lib}}^n(X)$}}

 b) $V$ opérant sur $H^i(X)\otimes_W k$ pour $i<n$ est nilpotent, et aussi opérant sur $\operatorname{Tors}^i(X)\otimes_W k$ et ${}_p\operatorname{Tors}^i(X)$ pour \uuline{$i\leq n$}. On a
\margin{$(**)$}
\diagram{\begin{tikzcd}[column sep=small]
0\arrow[r]& (H^n_{\mathrm{lib}}(X)\otimes k)_{\unsure{\mathrm{ss}}}^{\mathrm{ss}}\arrow[r]&H^0(X,\underline{\Omega}_X^n)^{\mathrm{ss}}\arrow[r,Rightarrow]&{}_p\operatorname{Tors}^{n+1}(X)_{\unsure{\mathrm{ss}}}^{\mathrm{ss}}\arrow[r]&0\unsure{\operatorname{Tor}}
\end{tikzcd}}

Sans doute $(*)$ et $(**)$ sont transposés l'un de l'autre \ldots

\uline{Notation}\quad L'indice ${}_{\mathrm{ss}}$ indique partie semi-simple pour $F$\\
L'exposant ${}^{\mathrm{ss}}$ \rule{7cm}{.3pt} $V$

\mspage{1}
